\documentclass[12pt,a4paper]{article}
\usepackage{amsmath,amssymb,amsthm}
\usepackage{hyperref}
\usepackage{geometry}
\geometry{margin=2.5cm}
\usepackage{enumitem}

\newtheorem{theorem}{Theorem}[section]
\newtheorem{lemma}[theorem]{Lemma}
\theoremstyle{definition}
\newtheorem{definition}[theorem]{Definition}
\theoremstyle{remark}
\newtheorem{remark}[theorem]{Remark}

\title{Ultra-Diffuse Galaxies from\\
Environment-Dependent Substrate Coherence}

\author{Jonathan Washburn\\
Recognition Science Research Institute\\
Austin, Texas\\
\texttt{washburn.jonathan@gmail.com}}

\date{March 2026}

\begin{document}
\maketitle

\begin{abstract}
We analyze ultra-diffuse galaxies (UDGs) within Recognition Science
(RS).  UDGs have effective radii comparable to the Milky Way
($r_e \gtrsim 1.5$~kpc) but luminosities $100$--$1000\times$ lower,
challenging hierarchical $\Lambda$CDM structure formation.  Some UDGs
(e.g., NGC~1052-DF2, DF4) appear to lack dark matter
entirely~\cite{vanDokkum2018,vanDokkum2019}, while others have
enormous dark-to-luminous ratios.  RS offers a qualitative explanation through Interacting Ledger
Gravity (ILG): dark matter is a substrate property whose effective
coupling depends on local substrate coherence $C_{\mathrm{sub}}$.
UDGs form in low-coherence regions of the cosmic web where the ILG
tidal coefficient $\alpha_t$ is suppressed, producing extended,
diffuse stellar distributions.  ``Dark-matter-free'' UDGs arise when
tidal interactions with neighboring massive halos locally deplete
substrate coherence.  Dark-matter-rich UDGs inhabit high-coherence
filament nodes.  We note that the ``dark-matter-free'' status of
NGC~1052-DF2 has been contested on the grounds of uncertain distance
measurement; the RS analysis is qualitative and accommodates a range
of observational outcomes.
\end{abstract}

%% ============================================================
\section{Introduction}
\label{sec:intro}
%% ============================================================

In the Recognition Science~\cite{RS_Axioms} (RS) framework, dark matter is
not a particle species but a property of the recognition substrate;
its effective strength is environment-dependent through the
Interacting Ledger Gravity (ILG) coupling.  This provides a natural
framework for UDG diversity.

Ultra-diffuse galaxies are a class of large, low-surface-brightness
systems first systematically identified in the Coma cluster by van
Dokkum et al.~\cite{vanDokkum2015}.  They present two puzzles for
$\Lambda$CDM:

\begin{enumerate}[nosep]
  \item \textbf{Large sizes with low luminosity.}  UDGs have
  $r_e \sim 1.5$--$5$~kpc (comparable to $L^\ast$ galaxies) but
  $M_V \gtrsim -15$ (dwarf-galaxy luminosity).  Standard CDM halo
  models struggle to produce such extended, low-mass systems.

  \item \textbf{Extreme dark matter fractions.}  Some UDGs appear
  ``dark-matter-free'' (NGC~1052-DF2: $M_{\mathrm{DM}}/M_\ast
  \lesssim 1$~\cite{vanDokkum2018}), while others have
  $M_{\mathrm{DM}}/M_\ast \sim 10^3$ (Dragonfly~44~\cite{vanDokkum2016}).
  This bimodality is difficult to explain with a single particle
  species.
\end{enumerate}

%% ============================================================
\section{RS Analysis: Interacting Ledger Gravity}
\label{sec:analysis}
%% ============================================================

\begin{definition}[Substrate Coherence]
The \emph{substrate coherence} $C_{\mathrm{sub}}$ at a given location
is the local density of aligned recognition events in the ledger.
$C_{\mathrm{sub}}$ varies with environment: high in filament nodes,
low in voids and tidal disruption zones.
\end{definition}

\begin{proposition}[UDGs from Low Substrate Coherence]
\label{thm:udg}
UDGs form in cosmic web regions with $C_{\mathrm{sub}} <
C_{\mathrm{crit}}$ (the critical coherence threshold for normal
galaxy formation).  In low-coherence regions:
\begin{enumerate}[label=(\roman*),nosep]
  \item The ILG tidal coefficient $\alpha_t$ is suppressed, weakening
  the effective gravitational coupling.
  \item Stars form in extended, diffuse distributions.
  \item Surface brightness is low because the stellar density is
  spread over a large volume.
\end{enumerate}
\end{proposition}

\begin{proof}[Structural argument]
In ILG, the effective gravitational coupling is proportional to the
local substrate coherence $C_{\mathrm{sub}}$.  When $C_{\mathrm{sub}} <
C_{\mathrm{crit}}$, the effective enclosed dark matter density is
reduced, lowering the restoring gravitational force on gas clouds during
collapse.  Stars form at lower density thresholds, producing extended
distributions.  This argument is qualitative; the quantitative value
of $C_{\mathrm{crit}}$ and the derivation of $\alpha_t(C_{\mathrm{sub}})$
from first principles are open problems in the ILG framework.
\end{proof}

\begin{proposition}[Dark-Matter-Free UDGs]
\label{thm:df2}
Tidal interactions with massive neighboring halos can locally deplete
substrate coherence, producing regions where
$C_{\mathrm{sub}} \approx 0$.  Galaxies forming in such regions
have negligible effective dark matter:
$M_{\mathrm{DM}}^{\mathrm{eff}} \approx 0$.
\end{proposition}

\begin{proof}[Structural argument]
In ILG, the effective dark matter $M_{\mathrm{DM}}^{\mathrm{eff}} \propto
C_{\mathrm{sub}}$.  Strong tidal fields from a neighboring massive halo
(such as NGC~1052 for DF2) can disrupt the local coherence of the
recognition substrate, driving $C_{\mathrm{sub}} \to 0$ in the stripped
region.  This provides a qualitative explanation for apparent
dark-matter depletion; a quantitative tidal stripping rate from the RS
J-cost framework has not been derived.
\end{proof}

\begin{proposition}[Dark-Matter-Rich UDGs]
\label{thm:df44}
UDGs at high-coherence filament nodes have enhanced effective dark
matter: $M_{\mathrm{DM}}^{\mathrm{eff}}/M_\ast \gg 1$.  The
substrate coherence amplifies the ILG effect, producing large apparent
dark matter fractions.
\end{proposition}

\begin{proof}[Structural argument]
Complementary to Proposition~\ref{thm:df2}: when $C_{\mathrm{sub}}$
is large (filament node), the ILG tidal enhancement amplifies the
effective dark matter density.  Dragonfly~44, residing in the Coma
cluster filament, exemplifies this regime.  The bimodality of UDG
dark matter fractions is a direct consequence of the environmental
dependence of $C_{\mathrm{sub}}$.
\end{proof}

\begin{remark}
In RS, dark matter is not a fixed mass associated with each galaxy
but an environment-dependent substrate property.  The same framework
that suppresses effective dark matter in tidal disruption zones
enhances it in coherence nodes.

\medskip\noindent\textbf{Note on NGC~1052-DF2.}  The ``dark-matter-free''
conclusion for NGC~1052-DF2 rests on the assumption that the galaxy
lies at the same distance as NGC~1052 ($\approx 20$~Mpc).  Several
authors have argued for a closer distance ($\approx 13$~Mpc), which
would reduce both the inferred size and the apparent dark matter deficit.
The RS analysis is consistent with either interpretation: low coherence
(dark matter deficit) or closer proximity (less extreme case).  The
observational situation for DF2 remains contested.
\end{remark}

%% ============================================================
\section{Predictions}
\label{sec:predict}
%% ============================================================

\begin{enumerate}
  \item \textbf{Environmental correlation.}  UDGs with low apparent
  dark matter fractions should preferentially be found near massive
  halos (in tidal disruption zones).  UDGs with high dark matter
  fractions should be found in filament nodes.  This is testable with
  large UDG surveys.

  \item \textbf{No dark matter particle.}  The dark matter content of
  UDGs is a substrate property, not a particle density.  No WIMP or
  axion detection experiment will find the dark matter ``missing''
  from DF2/DF4.

  \item \textbf{Transition population.}  Between the dark-matter-free
  and dark-matter-rich extremes, there should be a continuous
  distribution of UDG dark matter fractions, correlated with local
  substrate coherence.

  \item \textbf{Velocity dispersion.}  The stellar velocity dispersion
  of UDGs should correlate with local substrate coherence rather than
  with the galaxy's stellar mass alone.
\end{enumerate}

%% ============================================================
\section{Conclusion}
\label{sec:conclusion}
%% ============================================================

Ultra-diffuse galaxies are a natural consequence of the
environment-dependent substrate coherence in the ILG framework.
The bimodal distribution of dark matter fractions (from
dark-matter-free to dark-matter-dominated) arises from the spatial
variation of substrate coherence in the cosmic web.

%% ============================================================
\begin{thebibliography}{99}

\bibitem{vanDokkum2015}
P.~van~Dokkum et al., ApJ Lett.\ \textbf{798}, L45 (2015).

\bibitem{vanDokkum2016}
P.~van~Dokkum et al., ApJ Lett.\ \textbf{828}, L6 (2016).

\bibitem{vanDokkum2018}
P.~van~Dokkum et al., Nature \textbf{555}, 629 (2018).

\bibitem{vanDokkum2019}
P.~van~Dokkum et al., ApJ Lett.\ \textbf{874}, L5 (2019).

\bibitem{RS_Axioms}
J.~Washburn, ``Recognition Science: Derivation of the Cost
Functional and Forcing Chain,'' Recognition Science Research
Institute Technical Report (2025).

\end{thebibliography}

\end{document}
